% ==========================================
% BAB VI EVALUASI
% ==========================================
\cleardoublepage
\chapter{EVALUASI}
\label{chap:evaluasi}

\section{Tujuan Evaluasi}
\label{sec:tujuan-evaluasi}

Evaluasi pada tugas akhir ini bertujuan untuk menilai apakah rancangan
arsitektur sistem otomasi inspeksi kontainer yang telah direalisasikan
mendukung kebutuhan utama yang telah dirumuskan pada
Bab~\ref{chap:perancangan}. Berbeda dengan evaluasi yang berfokus pada akurasi
model pada kumpulan data uji operasional penuh, evaluasi pada tugas akhir ini menitikberatkan pada
kesesuaian realisasi arsitektur, integrasi antarkomponen, konsistensi alur data,
dan dukungan sistem terhadap proses inspeksi secara terpadu.

Pada evaluasi ini, kebutuhan utama didefinisikan sebagai kebutuhan berprioritas
tinggi pada Tabel~\ref{tbl:kebutuhan_utama_rancangan} yang langsung menentukan
alur inspeksi kontainer: integrasi komponen \textit{edge}, API, web, dan penyimpanan;
konsistensi identitas kontainer; pembentukan riwayat kontainer dan bukti visual;
pemisahan jalur data terstruktur dari jalur video langsung; serta kontrol akses
dan pencatatan aktivitas.

Secara khusus, evaluasi diarahkan untuk menjawab empat pertanyaan utama:

\begin{enumerate}
  \item Apakah komponen \textit{edge}, layanan API, aplikasi web, dan lapisan
  penyimpanan dapat berkomunikasi sesuai dengan rancangan arsitektur?
  \item Apakah alur data utama dari identifikasi kontainer sampai penyajian
  hasil dapat berjalan secara konsisten?
  \item Apakah struktur data yang direalisasikan mendukung riwayat proses
  inspeksi pada tingkat kontainer?
  \item Apakah pemisahan jalur data terstruktur dan jalur video langsung
  memenuhi indikator dukungan operasional pada ruang lingkup evaluasi?
\end{enumerate}

\section{Metode Evaluasi}
\label{sec:metode-evaluasi}

\subsection{Pendekatan Evaluasi}
\label{subsec:pendekatan-evaluasi}

Metode evaluasi yang digunakan adalah evaluasi berbasis skenario integrasi dan
validasi fungsional sistem. Pendekatan ini dipilih karena sesuai dengan tujuan
tugas akhir ini, yaitu menilai implementabilitas rancangan arsitektur dan kesesuaian
realisasinya terhadap kebutuhan sistem. Artefak utama yang digunakan sebagai
dasar evaluasi adalah dokumentasi pengujian integrasi, struktur layanan
yang telah direalisasikan, struktur skema data operasional, perilaku sistem yang telah direalisasikan pada
proyek pengembangan sistem, serta dokumen lapangan yang menunjukkan konteks
operasional \textit{screening} di lingkungan pelabuhan.

Karena fokus tugas akhir ini berada pada arsitektur sistem, evaluasi tidak
diarahkan untuk menilai kualitas rekayasa perangkat lunak pada tingkat kode,
gaya implementasi antarmuka, atau performa tiap model secara terpisah. Yang
dinilai adalah apakah artefak realisasi yang tersedia
merealisasikan keputusan arsitektural yang telah dirumuskan pada tahap
perancangan.

Evaluasi dilakukan terhadap tiga komponen utama sistem, yaitu komponen
\textit{edge} yang bertugas menarik aliran video, menjalankan inferensi, dan
menyiarkan hasil secara waktu nyata; layanan API yang menerima data hasil
inspeksi, mengelola persistensi, dan menyediakan antarmuka data terstruktur;
serta aplikasi web yang berfungsi sebagai antarmuka pemantauan, pengelolaan
aliran, dan pembacaan hasil inspeksi.

Selain itu, evaluasi juga mencakup komponen penyimpanan, yaitu MongoDB sebagai
basis data utama dan Cloudflare R2 sebagai penyimpanan artefak visual.

Dengan demikian, evaluasi pada bab ini membaca kembali matriks pemetaan
rancangan pada Tabel~\ref{tbl:traceability-rancangan-arsitektur}. Setiap
skenario evaluasi diposisikan sebagai pemeriksaan terhadap keputusan
arsitektural, bukan sebagai pengujian seluruh fitur aplikasi atau pengukuran
performansi skala operasional secara menyeluruh.

\subsection{Skenario Evaluasi}
\label{subsec:skenario-evaluasi}

Berdasarkan artefak pengujian integrasi yang tersedia, evaluasi dilakukan
melalui beberapa skenario utama.

\begin{enumerate}
  \item Verifikasi konektivitas \textit{edge} ke layanan API dilakukan untuk
  memastikan bahwa komponen \textit{edge} mengirimkan hasil identifikasi
  dan hasil pemindaian ke layanan API dengan mekanisme autentikasi yang sesuai.
  Hasil sesuai indikator pada skenario ini menunjukkan bahwa jalur data
  terstruktur dari lapisan \textit{edge} ke lapisan aplikasi berjalan dengan
  benar. Pada implementasi
  aktual, mekanisme ini direalisasikan melalui \textit{API key} pada keluarga
  \textit{endpoint} \texttt{/api/edge}.

  \item Verifikasi aliran OCR dilakukan untuk memeriksa apakah aliran OCR dapat
  dijalankan, mendeteksi nomor kontainer, dan menghasilkan entitas kontainer pada
  layanan aplikasi. Hasil sesuai indikator pada skenario ini menunjukkan bahwa
  identitas utama objek inspeksi dapat dibentuk secara otomatis dan tersedia untuk proses
  lanjutan. Dalam implementasi, proses ini diwujudkan melalui pembuatan atau
  pembaruan dokumen kontainer pada koleksi \textit{containers}.

  \item Verifikasi propagasi nomor kontainer digunakan untuk menilai apakah hasil
  deteksi nomor kontainer dari aliran OCR dapat dipropagasikan ke aliran lain,
  seperti aliran pemindaian kerusakan. Skenario ini penting karena
  menunjukkan bahwa integrasi lintas modul tidak bergantung pada pengisian manual
  yang berulang.

  \item Verifikasi penyimpanan hasil pemindaian dilakukan untuk memeriksa apakah
  hasil pemindaian kerusakan, barang berisiko, dan kategori muatan dapat
  tersimpan pada entitas kontainer yang sesuai. Skenario ini menjadi indikator
  bahwa model data yang dirancang digunakan sebagai wadah integrasi
  informasi inspeksi. Selain itu, skenario ini juga memeriksa bahwa hasil
  pemindaian tetap berada di dalam struktur data inti, bukan pada koleksi
  pelacakan tambahan.

  \item Verifikasi siaran video langsung dilakukan untuk menilai apakah aplikasi
  web dapat menerima aliran video beranotasi langsung dari komponen \textit{edge}
  melalui \textit{WebSocket}. Hasil sesuai indikator pada skenario ini
  menunjukkan bahwa jalur video langsung yang dirancang dapat direalisasikan
  tanpa membebani layanan API.

  \item Verifikasi perilaku antarmuka web dilakukan untuk memeriksa apakah
  aplikasi web dapat menampilkan aliran aktif, memperbarui status tab,
  menampilkan kontainer aktif hasil OCR, dan memperbarui halaman inspeksi setelah
  proses \textit{scan} diselesaikan. Evaluasi ini menunjukkan bahwa hasil
  integrasi antarkomponen dapat diterjemahkan menjadi interaksi pengguna yang
  utuh.
\end{enumerate}

Skenario tersebut dibaca dalam tiga kondisi evaluasi agar klaim hasil
tidak bersifat umum. Kondisi terbaik menunjukkan jalur yang diharapkan berjalan
lengkap, kondisi nominal menunjukkan alur yang masih berjalan dengan konfirmasi
operator, sedangkan kondisi terburuk menunjukkan batas ketika riwayat kontainer
tidak dapat dibentuk secara lengkap. Rangkuman kondisi evaluasi ditunjukkan pada
Tabel~\ref{tbl:kondisi-skenario-evaluasi}.

\begin{table}[ht]
  \centering
  \scriptsize
  \caption[Kondisi skenario evaluasi]{Kondisi Skenario Evaluasi}
  \label{tbl:kondisi-skenario-evaluasi}
  \setlength{\tabcolsep}{3pt}
  \renewcommand{\arraystretch}{1.12}
  \begin{tabularx}{\textwidth}{>{\raggedright\arraybackslash}p{0.17\textwidth} >{\raggedright\arraybackslash}X >{\raggedright\arraybackslash}X >{\raggedright\arraybackslash}X}
    \toprule
    \textbf{Kondisi} & \textbf{Pemicu} & \textbf{Bukti yang Diperiksa} & \textbf{Batas Klaim} \\
    \midrule
    Kondisi terbaik (\textit{best case}) & Koneksi \textit{edge}--API tersedia, nomor kontainer terbaca OCR, pemindaian lain aktif, artefak visual tersimpan, dan aplikasi web menerima video langsung. & Data OCR dan hasil pemindaian muncul pada entitas kontainer yang sama, bukti visual tersedia, dan halaman web dapat membaca hasilnya. & Rancangan ditunjukkan berjalan pada alur internal yang memiliki identitas, hasil pemindaian, dan bukti visual. \\
    Kondisi nominal & Nomor kontainer terbaca tetapi sebagian hasil membutuhkan verifikasi operator atau artefak visual baru tersedia setelah proses pemindaian selesai. & Data kontainer tetap menjadi jangkar agregasi, status inspeksi dapat diperbarui, dan hasil dapat dibaca kembali melalui web. & Rancangan mendukung koreksi/konfirmasi manusia tanpa mengubah pusat data kontainer. \\
    Kondisi terburuk (\textit{worst case}) & Nomor kontainer tidak tersedia, koneksi API terganggu, atau artefak visual gagal tersimpan. & Ketidaklengkapan riwayat dapat diidentifikasi dari tidak adanya nomor kontainer, status pengiriman, atau referensi artefak visual. & Rancangan belum diklaim membentuk riwayat lengkap ketika identitas atau bukti utama hilang. \\
    \bottomrule
  \end{tabularx}
\end{table}

\FloatBarrier

\subsection{Estimasi Kapasitas Inspeksi}
\label{subsec:estimasi-kapasitas-inspeksi}

Selain skenario integrasi, rancangan juga perlu dibaca dari sisi kapasitas
operasional. Karena evaluasi tugas akhir ini belum berupa uji lapangan penuh,
kapasitas dihitung sebagai estimasi berbasis skenario. Estimasi ini digunakan
untuk menunjukkan bagaimana arsitektur dapat diskalakan ketika jumlah perangkat
\textit{edge} atau jalur inspeksi ditambah.

Kapasitas inspeksi per jam dihitung dengan Persamaan~\ref{eq:kapasitas-inspeksi}.

\begin{equation}
K = \frac{3600}{t_s} \times n \times u
\label{eq:kapasitas-inspeksi}
\eqcaption{Estimasi kapasitas inspeksi per jam}
\end{equation}

dengan \(K\) adalah estimasi jumlah kontainer per jam, \(t_s\) adalah waktu
siklus inspeksi per kontainer dalam detik, \(n\) adalah jumlah jalur atau
perangkat inspeksi yang berjalan paralel, dan \(u\) adalah faktor utilisasi
operasional. Faktor utilisasi digunakan karena tidak seluruh waktu operasional
dapat dipakai untuk pemrosesan aktif; terdapat waktu tunggu, reposisi
kontainer, konfirmasi operator, atau gangguan jaringan.

Nilai \(t_s\) pada estimasi ini diperlakukan sebagai asumsi skenario, bukan
hasil pengukuran lapangan penuh. Nilai 30 detik/kontainer merepresentasikan
alur ketika kamera, OCR, pemindaian visual, dan pengiriman hasil berjalan tanpa
hambatan berarti. Nilai 60 detik/kontainer digunakan sebagai kondisi nominal
ketika terdapat waktu konfirmasi operator atau penyesuaian posisi kontainer.
Nilai 120 detik/kontainer digunakan sebagai kondisi konservatif ketika terjadi
pengulangan pembacaan, gangguan jaringan sementara, atau kebutuhan validasi
manual. Dengan demikian, angka kapasitas pada Tabel~\ref{tbl:estimasi-kapasitas-inspeksi}
dibaca sebagai batas rancangan untuk membandingkan skenario, bukan sebagai
klaim kapasitas operasional final pelabuhan.

\begin{table}[ht]
  \centering
  \small
  \caption[Estimasi kapasitas inspeksi]{Estimasi Kapasitas Inspeksi Berdasarkan Skenario}
  \label{tbl:estimasi-kapasitas-inspeksi}
  \setlength{\tabcolsep}{4pt}
  \renewcommand{\arraystretch}{1.1}
  \begin{tabularx}{\textwidth}{>{\raggedright\arraybackslash}p{0.18\textwidth} >{\raggedright\arraybackslash}p{0.20\textwidth} >{\raggedright\arraybackslash}p{0.14\textwidth} >{\raggedright\arraybackslash}p{0.18\textwidth} >{\raggedright\arraybackslash}X}
    \toprule
    \textbf{Skenario} & \textbf{Asumsi Waktu Siklus} & \textbf{Jumlah Jalur} & \textbf{Utilisasi} & \textbf{Estimasi Kapasitas} \\
    \midrule
    Kondisi terbaik & 30 detik/kontainer & 1 & 0,80 & 96 kontainer/jam \\
    Kondisi nominal & 60 detik/kontainer & 1 & 0,70 & 42 kontainer/jam \\
    Kondisi terburuk & 120 detik/kontainer & 1 & 0,50 & 15 kontainer/jam \\
    Skalasi dua jalur & 60 detik/kontainer & 2 & 0,70 & 84 kontainer/jam \\
    Skalasi empat jalur & 60 detik/kontainer & 4 & 0,70 & 168 kontainer/jam \\
    \bottomrule
  \end{tabularx}
\end{table}

\FloatBarrier

Perhitungan pada Tabel~\ref{tbl:estimasi-kapasitas-inspeksi} tidak dimaksudkan
sebagai hasil uji performa akhir, melainkan sebagai batas rancangan yang dapat
ditelusuri. Nilai \(t_s\) mencakup akuisisi citra, pembacaan nomor kontainer,
pemindaian visual, pengiriman hasil ke layanan API, penyimpanan data, dan
konfirmasi operator apabila diperlukan. Karena komponen \textit{edge} dapat
ditempatkan per jalur inspeksi, kapasitas rancangan meningkat secara hampir
linear terhadap jumlah jalur selama layanan API, basis data, penyimpanan objek,
dan jaringan masih memenuhi prasyarat beban untuk menerima tambahan data dari
jalur tersebut. Dengan demikian, estimasi ini
mendukung keputusan arsitektural untuk memisahkan pemrosesan lokal dari
persistensi terpusat.

\subsection{Aspek yang Dievaluasi}
\label{subsec:aspek-dievaluasi}

Berdasarkan skenario evaluasi pada Subbab~\ref{subsec:skenario-evaluasi}, aspek
evaluasi pada tugas akhir ini dirangkum dalam Tabel~\ref{tbl:aspek-evaluasi}.

\begin{table}[ht]
    \centering
    \small
    \caption[Aspek evaluasi arsitektur]{Aspek Evaluasi Realisasi Arsitektur}
    \label{tbl:aspek-evaluasi}
    \begin{tabularx}{\textwidth}{>{\raggedright\arraybackslash}p{0.24\textwidth} >{\raggedright\arraybackslash}p{0.25\textwidth} >{\raggedright\arraybackslash}X}
        \toprule
        \textbf{Aspek} & \textbf{Fokus Pengujian} & \textbf{Indikator Evaluasi} \\
        \midrule
        Integrasi layanan & Komunikasi \textit{edge}, API, dan web & Data dapat dikirim, diterima, dan dibaca antarkomponen \\
        Konsistensi alur data & Perpindahan data dari identifikasi hingga hasil inspeksi & Data hasil inspeksi terhubung ke entitas kontainer yang benar \\
        Riwayat kontainer & Pencatatan status, waktu, dan artefak visual & Tersedia jejak proses dan bukti inspeksi yang terhubung ke identitas kontainer \\
        Pemisahan jalur komunikasi & Pemisahan jalur data terstruktur dan video & Data transaksional dan video menggunakan jalur yang berbeda sesuai fungsi \\
        Kesiapan operasional antarmuka & Perilaku aplikasi web dalam menampilkan hasil & Operator dapat memantau aliran dan membaca hasil inspeksi secara utuh \\
        \bottomrule
    \end{tabularx}
\end{table}

\FloatBarrier

\textbf{Kriteria keterpenuhan evaluasi.}
Untuk menghindari istilah yang tidak terukur, status pada evaluasi ini dibaca
berdasarkan makna operasional pada Tabel~\ref{tbl:kriteria-keterpenuhan-evaluasi}.
Dengan demikian, kata seperti ``terpenuhi'', ``memadai'', dan ``lengkap'' tidak
digunakan sebagai penilaian umum, tetapi sebagai status yang bergantung pada
indikator dan bukti yang diperiksa.

\begin{table}[ht]
    \centering
    \scriptsize
    \caption[Kriteria keterpenuhan evaluasi]{Kriteria Keterpenuhan Evaluasi}
    \label{tbl:kriteria-keterpenuhan-evaluasi}
    \setlength{\tabcolsep}{4pt}
    \renewcommand{\arraystretch}{1.12}
    \begin{tabularx}{\textwidth}{>{\raggedright\arraybackslash}p{0.18\textwidth} >{\raggedright\arraybackslash}X >{\raggedright\arraybackslash}X >{\raggedright\arraybackslash}X}
        \toprule
        \textbf{Status/Istilah} & \textbf{Makna Operasional} & \textbf{Indikator Minimal} & \textbf{Batas Klaim} \\
        \midrule
        Terpenuhi internal & Keputusan arsitektural muncul pada artefak realisasi internal yang diuji. & Ada bukti struktur data, alur API, tampilan web, atau artefak penyimpanan yang menunjukkan perilaku tersebut. & Tidak berarti sudah tervalidasi pada operasi pelabuhan penuh. \\
        Terpenuhi terbatas & Bukti tersedia, tetapi hanya pada skenario, data, atau artefak tertentu. & Minimal satu skenario menunjukkan hasil yang sesuai indikator, tetapi variasi kondisi belum diuji menyeluruh. & Tidak digunakan untuk mengklaim akurasi, performa, atau reliabilitas umum. \\
        Terdukung arsitektural & Rancangan menyediakan mekanisme, kontrak, atau pemisahan komponen untuk memenuhi kebutuhan. & Ada desain komponen, model data, atau pola integrasi yang dapat ditelusuri ke kebutuhan. & Belum sama dengan realisasi operasional penuh. \\
        Memadai & Indikator pada Tabel~\ref{tbl:aspek-evaluasi} terpenuhi untuk ruang lingkup skenario yang diuji. & Bukti yang diperiksa cukup untuk menjawab pertanyaan evaluasi pada ruang lingkup tugas akhir. & Tidak berarti cukup untuk sertifikasi, audit eksternal, atau beban produksi. \\
        Lengkap & Riwayat kontainer memiliki identitas kontainer, waktu/status proses, hasil inspeksi, dan referensi artefak visual. & Keempat unsur tersebut muncul pada entitas kontainer atau artefak evaluasi. & Tidak diklaim lengkap jika identitas atau bukti visual utama hilang. \\
        \bottomrule
    \end{tabularx}
\end{table}

\FloatBarrier

\subsection{Matriks Bukti Evaluasi}
\label{subsec:matriks-bukti-evaluasi}

Tabel~\ref{tbl:aspek-evaluasi} berfungsi sebagai rubrik aspek yang dinilai,
sedangkan hasil validasinya ditunjukkan secara lebih konkret pada
Tabel~\ref{tbl:matriks-bukti-evaluasi}. Matriks tersebut memperlihatkan cara
validasi, bukti yang diperiksa, hasil aktual yang diamati, dan status
keterpenuhan setiap skenario. Dengan demikian, evaluasi tidak hanya berupa
pernyataan umum tentang hasil rancangan, tetapi menunjukkan dasar
pemeriksaannya.

Prosedur validasi arsitektur dilakukan dalam empat langkah. Pertama, setiap
skenario dikaitkan dengan keputusan arsitektural yang telah dirumuskan pada
Bab~\ref{chap:perancangan}, seperti pemisahan komponen, penggunaan entitas
kontainer sebagai pusat agregasi, pemisahan jalur video dari jalur data
terstruktur, dan pemisahan penyimpanan artefak visual. Kedua, artefak realisasi
yang diperlukan diperiksa, meliputi struktur layanan API, struktur data kontainer,
tampilan antarmuka, dan bukti penyimpanan artefak visual. Ketiga, hasil aktual
dibandingkan dengan indikator pada Tabel~\ref{tbl:aspek-evaluasi}.
Keempat, skenario diberi status sesuai kriteria pada
Tabel~\ref{tbl:kriteria-keterpenuhan-evaluasi}. Status pada tabel tidak
dimaksudkan sebagai klaim bahwa sistem telah divalidasi pada skala operasional
atau lingkungan operasional penuh.

\begin{table}[ht]
    \centering
    \scriptsize
    \caption[Matriks hasil validasi arsitektur]{Matriks Hasil Validasi Realisasi Arsitektur}
    \label{tbl:matriks-bukti-evaluasi}
    \setlength{\tabcolsep}{3pt}
    \renewcommand{\arraystretch}{1.08}
    \begin{tabularx}{\textwidth}{>{\raggedright\arraybackslash}p{0.18\textwidth} >{\raggedright\arraybackslash}X >{\raggedright\arraybackslash}X >{\raggedright\arraybackslash}X >{\raggedright\arraybackslash}p{0.13\textwidth}}
        \toprule
        \textbf{Skenario} & \textbf{Cara Validasi} & \textbf{Bukti Konkret} & \textbf{Hasil Aktual} & \textbf{Status} \\
        \midrule
        Konektivitas \textit{edge} ke layanan API &
        Memeriksa alur pengiriman hasil dari komponen \textit{edge} melalui keluarga \textit{endpoint} \texttt{/api/edge} dan mekanisme \textit{API key}. &
        Struktur layanan API dan bukti data pada entitas kontainer di Listing~\ref{lst:container-entity-sample}. &
        Hasil OCR, \textit{defect scan}, \textit{illegal scan}, dan \textit{category scan} tersimpan pada dokumen kontainer yang sama. &
        Terpenuhi internal \\
        OCR membentuk entitas kontainer &
        Memeriksa apakah hasil OCR menghasilkan atau memperbarui entitas kontainer sebagai pusat agregasi data. &
        Gambar~\ref{fig:evidence_ocr_live} dan Listing~\ref{lst:container-entity-sample}. &
        Nomor kontainer terdeteksi sebagai DFSU6625920 dan muncul kembali pada struktur data inti kontainer. &
        Terpenuhi terbatas \\
        Propagasi nomor kontainer ke alur lain &
        Memeriksa apakah nomor kontainer hasil OCR digunakan oleh alur pemindaian lain sebagai konteks bersama. &
        Gambar~\ref{fig:evidence_container_detail} dan struktur \texttt{defectScan}, \texttt{illegalScan}, serta \texttt{categoryScan} pada Listing~\ref{lst:container-entity-sample}. &
        Hasil pemindaian kerusakan, barang berisiko, dan kategori muatan dikaitkan ke entitas kontainer yang sama. &
        Terpenuhi terbatas \\
        Penyimpanan hasil inspeksi &
        Memeriksa lokasi penyimpanan data terstruktur dan artefak visual yang dihasilkan proses inspeksi. &
        Listing~\ref{lst:container-entity-sample} dan Gambar~\ref{fig:evidence_r2_storage}. &
        Data transaksional tersimpan di MongoDB, sedangkan bukti visual tersimpan sebagai objek pada Cloudflare R2. &
        Terpenuhi internal \\
        Siaran video langsung ke web &
        Memeriksa apakah aplikasi web menerima aliran video beranotasi dari komponen \textit{edge} melalui \textit{WebSocket}. &
        Gambar~\ref{fig:evidence_ocr_live} dan Gambar~\ref{fig:evidence_defect_live}. &
        Antarmuka web menampilkan aliran OCR dan \textit{defect scan} secara langsung dengan anotasi hasil deteksi. &
        Terpenuhi terbatas \\
        Pembacaan hasil oleh aplikasi web &
        Memeriksa apakah aplikasi web dapat membaca hasil inspeksi yang telah tersimpan melalui layanan API. &
        Gambar~\ref{fig:evidence_container_detail}. &
        Halaman detail kontainer menampilkan ringkasan OCR, jumlah \textit{defect}, \textit{illegal items}, kategori, jumlah \textit{scan}, dan \textit{metadata} inspeksi. &
        Terpenuhi internal \\
        \bottomrule
    \end{tabularx}
\end{table}

\FloatBarrier

\subsection{Keterkaitan Hasil Evaluasi dengan Kebutuhan Sistem}
\label{subsec:keterkaitan-hasil-kebutuhan}

Kebutuhan fungsional dan nonfungsional yang dirumuskan pada Bab~\ref{chap:analisis-masalah}
tidak seluruhnya dievaluasi sebagai pengujian produk akhir. Sebagian kebutuhan,
terutama yang berkaitan dengan integrasi eksternal dan tata kelola operasional
penuh, berada pada tingkat kesiapan rancangan karena ruang lingkup tugas akhir
ini adalah perancangan arsitektur sistem. Oleh karena itu, Tabel~\ref{tbl:pemetaan-fr-nfr-evaluasi}
merangkum hubungan antara kelompok kebutuhan, bukti evaluasi yang digunakan,
dan status validasinya.

\begin{table}[ht]
    \centering
    \scriptsize
    \caption[Pemetaan kebutuhan terhadap hasil evaluasi]{Pemetaan Kebutuhan terhadap Hasil Evaluasi Arsitektur}
    \label{tbl:pemetaan-fr-nfr-evaluasi}
    \setlength{\tabcolsep}{4pt}
    \renewcommand{\arraystretch}{1.12}
    \begin{tabularx}{\textwidth}{>{\raggedright\arraybackslash}p{0.18\textwidth} >{\raggedright\arraybackslash}p{0.18\textwidth} >{\raggedright\arraybackslash}X >{\raggedright\arraybackslash}p{0.18\textwidth}}
        \toprule
        \textbf{Kelompok Kebutuhan} & \textbf{Kode Terkait} & \textbf{Bukti Evaluasi} & \textbf{Status Validasi} \\
        \midrule
        Pencatatan riwayat inspeksi &
        FR-01, FR-02, FR-03, FR-10 &
        Entitas kontainer menyimpan hasil OCR, hasil pemindaian, status, waktu, penanggung jawab, dan referensi artefak visual sebagaimana ditunjukkan pada Listing~\ref{lst:container-entity-sample}. &
        Terpenuhi internal \\
        Standardisasi proses dan konsistensi interpretasi &
        FR-04, FR-05, FR-06 &
        Struktur data hasil inspeksi memisahkan jenis temuan, tingkat risiko, tingkat keparahan, dan konteks hasil pemindaian sehingga hasil dapat dibaca secara konsisten oleh aplikasi web. &
        Terpenuhi terbatas \\
        Integrasi data dan penyajian informasi &
        FR-07, FR-08, FR-09 &
        Keluarga \textit{endpoint} \nolinkurl{/api/edge}, antarmuka web, dan pola \textit{adapter} menunjukkan kesiapan pertukaran data. Integrasi penuh dengan TOS, PCS, atau sistem eksternal lain tidak diklaim sebagai hasil realisasi. &
        Terpenuhi internal untuk integrasi internal; integrasi eksternal tidak diklaim \\
        Kontrol akses dan pemisahan konteks penggunaan &
        FR-11, NFR-04 &
        Komponen \textit{edge} menggunakan \textit{API key}, sedangkan aplikasi web menggunakan autentikasi berbasis pengguna. &
        Terpenuhi terbatas \\
        Pemisahan jalur komunikasi dan pemantauan waktu nyata &
        NFR-01, NFR-03 &
        Video beranotasi disalurkan melalui \textit{WebSocket}, sedangkan data inspeksi dikirim melalui layanan API dan disimpan secara terstruktur. &
        Terdukung arsitektural; belum diuji sebagai performa kuantitatif menyeluruh \\
        Interoperabilitas bertahap &
        NFR-02 &
        Pola \textit{adapter}, kontrak data, dan pemisahan layanan inti menunjukkan bahwa integrasi dengan sistem eksternal dapat dikembangkan tanpa mengubah logika inspeksi utama. &
        Terdukung arsitektural sebagai kesiapan integrasi \\
        Kemudahan pemeliharaan &
        NFR-05 &
        Pemisahan komponen \textit{edge}, API, web, basis data, dan penyimpanan objek menunjukkan batas tanggung jawab yang jelas dan mendukung perubahan bertahap per komponen. &
        Terdukung arsitektural \\
        \bottomrule
    \end{tabularx}
\end{table}

\FloatBarrier

\section{Hasil Evaluasi}
\label{sec:hasil-evaluasi}

Sesuai arahan evaluasi pada tugas akhir ini, hasil tidak hanya dibaca sebagai
daftar fungsi yang berjalan, tetapi sebagai indikasi bahwa rancangan arsitektur
dapat direalisasikan pada skenario yang diuji dan mendukung tujuan yang telah
ditetapkan dalam ruang lingkup tugas akhir.
Karena itu, setiap hasil pada bagian ini perlu dipahami melalui tiga lapis:

\begin{enumerate}
  \item Premis kebutuhan operasional dan prinsip arsitektur yang menjadi dasar
  rancangan, seperti kebutuhan integrasi, riwayat kontainer, dan pemisahan tanggung
  jawab antarkomponen.
  \item Keputusan arsitektural yang diambil untuk menjawab kebutuhan tersebut.
  \item Artefak implementasi dan skenario evaluasi yang menunjukkan bahwa
  keputusan tersebut dapat direalisasikan pada sistem yang dibangun.
\end{enumerate}

\begin{table}[ht]
    \centering
    \scriptsize
    \caption[Ringkasan hasil evaluasi dan bukti operasional]{Ringkasan Hasil Evaluasi dan Bukti Operasional}
    \label{tbl:ringkasan-hasil-evaluasi}
    \setlength{\tabcolsep}{4pt}
    \renewcommand{\arraystretch}{1.12}
    \begin{tabularx}{\textwidth}{>{\raggedright\arraybackslash}p{0.20\textwidth} >{\raggedright\arraybackslash}X >{\raggedright\arraybackslash}X >{\raggedright\arraybackslash}X}
        \toprule
        \textbf{Hasil yang Dinilai} & \textbf{Indikator} & \textbf{Bukti yang Digunakan} & \textbf{Batas Klaim} \\
        \midrule
        Identitas kontainer terbentuk dari OCR & Nomor kontainer muncul pada aliran OCR dan menjadi identitas pada data kontainer. & Gambar~\ref{fig:evidence_ocr_live} dan Listing~\ref{lst:container-entity-sample}. & Terpenuhi terbatas pada artefak evaluasi yang tersedia. \\
        Hasil inspeksi terikat pada kontainer yang sama & Hasil OCR, \textit{defect scan}, \textit{illegal scan}, dan \textit{category scan} berada pada satu entitas kontainer. & Listing~\ref{lst:container-entity-sample} dan Gambar~\ref{fig:evidence_container_detail}. & Tidak mengklaim seluruh variasi kasus lapangan. \\
        Jalur video dan data terstruktur terpisah & Video beranotasi disalurkan melalui \textit{WebSocket}, sedangkan data hasil inspeksi disimpan melalui layanan API. & Gambar~\ref{fig:evidence_ocr_live}, Gambar~\ref{fig:evidence_defect_live}, dan matriks pada Tabel~\ref{tbl:matriks-bukti-evaluasi}. & Menunjukkan realisasi internal, bukan uji beban video skala penuh. \\
        Artefak visual tersimpan terpisah dari data inti & Data inti berada pada MongoDB, sedangkan bukti visual berada pada penyimpanan objek. & Listing~\ref{lst:container-entity-sample} dan Gambar~\ref{fig:evidence_r2_storage}. & Tidak menilai kebijakan retensi atau audit eksternal. \\
        Kapasitas inspeksi dapat diekstrapolasi & Kapasitas dihitung dari waktu siklus, jumlah jalur, dan utilisasi. & Persamaan kapasitas dan Tabel~\ref{tbl:estimasi-kapasitas-inspeksi}. & Estimasi rancangan, bukan hasil uji lapangan penuh. \\
        \bottomrule
    \end{tabularx}
\end{table}

\FloatBarrier

Evaluasi integrasi antarlayanan menunjukkan bahwa komponen \textit{edge},
layanan API, dan aplikasi web telah direalisasikan sesuai dengan pembagian
tanggung jawab yang dirancang. Komponen \textit{edge} dapat mengirimkan hasil
identifikasi dan hasil pemindaian ke layanan API melalui HTTP, sedangkan
aplikasi web dapat membaca kembali data yang telah disimpan melalui antarmuka
REST.

Pengujian juga menunjukkan bahwa autentikasi antarkomponen telah dipisahkan
sesuai konteks penggunaannya. Komponen \textit{edge} menggunakan mekanisme
\textit{API key} untuk mengirimkan data mesin-ke-mesin, sedangkan aplikasi web
menggunakan autentikasi berbasis pengguna. Pemisahan ini menunjukkan bahwa
arsitektur tidak hanya memisahkan jalur data, tetapi juga memisahkan konteks
akses sesuai dengan peran masing-masing komponen.

Pada alur identifikasi dan pemindaian, sistem membuat atau memperbarui entitas
kontainer berdasarkan nomor kontainer yang dihasilkan oleh proses OCR. Setelah
identitas kontainer terbentuk, aliran pemindaian lain dapat menggunakan
identitas tersebut sebagai konteks bersama untuk menyimpan hasil inspeksi. Hal
ini memperlihatkan
bahwa kontainer menjadi entitas agregasi utama pada arsitektur data
yang dibangun.

Pengujian propagasi nomor kontainer menunjukkan bahwa aliran OCR dapat
menyebarkan nomor kontainer yang terdeteksi ke aliran lain yang sedang
aktif. Dengan mekanisme ini, hasil pemindaian kerusakan, barang berisiko, dan
kategori muatan dapat langsung dikaitkan ke kontainer yang sama. Hasil tersebut
menunjukkan bahwa integrasi lintas aliran tidak lagi mengandalkan entri manual
berulang.

\Needspace{7\baselineskip}
Pada jalur video langsung, hasil evaluasi memperlihatkan bahwa aplikasi web
dapat menerima aliran video beranotasi langsung dari komponen \textit{edge}
melalui \textit{WebSocket}. Hal ini menunjukkan bahwa keputusan arsitektural
untuk memisahkan jalur video dari
jalur data terstruktur dapat direalisasikan pada skenario evaluasi.

Pemisahan ini memberikan dua keuntungan utama:

\begin{enumerate}
  \item Aplikasi web tetap dapat menampilkan kondisi lapangan secara waktu nyata.
  \item Layanan API tidak perlu bertindak sebagai perantara untuk data video
  berukuran besar.
\end{enumerate}

Dengan demikian, jalur video langsung menunjukkan dukungan terhadap kebutuhan
pemantauan pada ruang lingkup skenario evaluasi tanpa mengganggu jalur
penyimpanan data transaksional.

\subsection{Evaluasi Arsitektur Data}
\label{subsec:hasil-arsitektur-data}

Evaluasi terhadap struktur data menunjukkan bahwa model data yang direalisasikan
telah mendukung kebutuhan integrasi dan riwayat kontainer. Entitas \textit{Container}
menyimpan hasil OCR, ringkasan manifest, hasil pemindaian kerusakan,
barang berisiko, klasifikasi muatan, dan bukti inspeksi dalam satu entitas
terpadu. Sementara itu, data manifes menyediakan konteks deklarasi muatan yang
dapat dibandingkan dengan hasil deteksi aktual, sehingga keterkaitan antara
identitas kontainer, isi muatan, dan hasil inspeksi dapat dijaga secara
konsisten. Struktur ini juga memperjelas bahwa operasi inti sistem bertumpu pada
koleksi \textit{containers}, \textit{manifests}, dan \textit{users}, sedangkan
kanal pelacakan tambahan tidak menjadi dasar utama agregasi hasil inspeksi.

Pemisahan penyimpanan antara MongoDB dan Cloudflare R2 juga menunjukkan bahwa
implementasi mengikuti kebutuhan data yang berbeda. Data terstruktur disimpan
untuk mendukung kueri dan pembaruan status, sedangkan artefak visual disimpan
di media penyimpanan objek yang lebih sesuai. Hasil ini menguatkan bahwa
rancangan arsitektur data yang diusulkan bersifat operasional, bukan sekadar
konseptual.

\subsection{Temuan Evaluasi Utama}
\label{subsec:temuan-evaluasi-utama}

Dari seluruh skenario dan artefak evaluasi, terdapat tiga temuan utama yang
paling kuat dalam mendukung kontribusi tugas akhir ini:

\begin{enumerate}
  \item Integrasi antarlayanan dapat direalisasikan pada jalur yang memang
  kritis bagi operasi sistem. Dokumentasi pengujian integrasi menunjukkan bahwa
  komponen \textit{edge}, layanan API, dan aplikasi web telah terhubung melalui
  pembagian jalur yang jelas: data terstruktur dikirim ke layanan API, sedangkan
  video langsung dikirim ke aplikasi web melalui \textit{WebSocket}. Temuan ini
  penting karena menunjukkan bahwa keputusan arsitektural inti tidak berhenti
  pada tingkat diagram, tetapi tercermin dalam perilaku sistem.
  Contoh hasil realisasi tersebut dapat dilihat pada
  Gambar~\ref{fig:evidence_container_detail} yang menunjukkan bahwa satu halaman
  detail kontainer dapat memuat ringkasan hasil OCR, jumlah \textit{defect},
  jumlah \textit{illegal items}, kategori muatan, jumlah \textit{scan}, dan
  \textit{metadata} inspeksi dalam satu tampilan teragregasi.

  \item Nomor kontainer berfungsi sebagai jangkar integrasi data lintas proses
  inspeksi. Hasil evaluasi pada alur OCR dan propagasi nomor kontainer
  memperlihatkan bahwa identitas kontainer tidak hanya dikenali, tetapi juga
  digunakan sebagai konteks bersama untuk mengaitkan hasil pemindaian lain.
  Dengan demikian, struktur data yang dirancang mendukung agregasi
  hasil inspeksi dari beberapa aliran ke satu entitas yang sama. Temuan ini
  menjadi bukti paling langsung bahwa fokus arsitektur data pada tugas akhir ini
  memiliki dasar evaluasi. Gambar~\ref{fig:evidence_ocr_live} menunjukkan bahwa nomor
  kontainer terdeteksi pada aliran OCR, sedangkan
  Gambar~\ref{fig:evidence_container_detail} menunjukkan bahwa identitas yang
  sama muncul kembali pada halaman detail kontainer sebagai pusat agregasi hasil
  inspeksi.

  \item Kebutuhan riwayat kontainer dan keandalan operasional didukung oleh konteks
  lapangan yang nyata. Dokumen lapangan terkait assessment alat pemindai,
  pengaturan lalu lintas alat pemindai, keberlangsungan layanan, dan prosedur
  pengelolaan insiden menunjukkan bahwa layanan \textit{screening} berada dalam
  lingkungan operasional yang menuntut koordinasi lintas pihak, kejelasan alur,
  dan kesiapan penanganan gangguan. Karena itu, keputusan arsitektural seperti
  pencatatan status layanan, jejak audit, dan pemisahan jalur komunikasi
  memiliki dasar kebutuhan operasional yang nyata, bukan sekadar preferensi
  desain. Bukti visual untuk jalur pemantauan langsung juga dapat diamati pada
  Gambar~\ref{fig:evidence_ocr_live} dan Gambar~\ref{fig:evidence_defect_live},
  yang memperlihatkan bahwa aliran OCR dan \textit{defect scan}
  tampil pada antarmuka pemantauan.
\end{enumerate}

\section{Pembahasan Hasil Evaluasi}
\label{sec:pembahasan-hasil-evaluasi}

Secara umum, hasil evaluasi menunjukkan bahwa rancangan arsitektur dapat
direalisasikan menjadi sistem yang terintegrasi pada ruang lingkup internal
tugas akhir. Dari perspektif arsitektur, nilai utama evaluasi tidak hanya
	terletak pada tersedianya komponen-komponen sistem, tetapi pada perilaku
	komponen tersebut ketika bekerja dalam alur yang konsisten.
Komponen \textit{edge} bertanggung jawab terhadap pengambilan data dan inferensi
awal, layanan API bertanggung jawab terhadap persistensi dan integrasi data,
sedangkan aplikasi web bertanggung jawab terhadap penyajian hasil dan interaksi
pengguna.

Hasil evaluasi juga memperlihatkan bahwa keputusan pemisahan jalur komunikasi
merupakan keputusan yang tepat. Jalur data terstruktur melalui layanan API
menjaga konsistensi dan kontrol akses, sedangkan jalur video langsung melalui
\textit{WebSocket} memenuhi kebutuhan pemantauan waktu nyata. Dengan kata lain,
arsitektur yang direalisasikan telah membedakan dengan jelas antara data yang
bersifat persisten dan data yang bersifat temporer.

Jika dikaitkan dengan konteks lapangan, pemisahan ini juga dibutuhkan secara
operasional. Layanan \textit{screening} dan inspeksi di pelabuhan melibatkan alur kerja
yang harus tetap bergerak, sementara bukti inspeksi dan status layanan perlu
tetap dapat ditelusuri. Karena itu, arsitektur tidak cukup hanya menyediakan
deteksi otomatis, tetapi juga harus menjaga agar data operasional, bukti
visual, dan status layanan dapat dikelola sesuai fungsi masing-masing.

Dari sisi kontribusi tugas akhir, hasil evaluasi mendukung argumentasi bahwa
fokus pada arsitektur sistem dan data memiliki dasar evaluasi. Struktur data yang
dibangun memungkinkan hasil inspeksi dari beberapa sumber untuk digabungkan
dalam satu entitas kontainer. Mekanisme ini penting dalam konteks pelabuhan
karena proses inspeksi jarang terjadi dalam satu titik tunggal. Sistem harus
menyatukan hasil dari beberapa tahapan inspeksi tanpa kehilangan konteks
dan tanpa mengorbankan riwayat kontainer. Oleh sebab itu, temuan evaluasi pada
tugas akhir ini perlu dibaca sebagai pembuktian terhadap rancangan arsitektur,
bukan sebagai penilaian menyeluruh terhadap seluruh keputusan rekayasa
perangkat lunak.

\subsection{Validasi Konseptual terhadap Tujuan Arsitektural}
\label{subsec:validasi-konseptual}

Jika dirumuskan dalam bentuk penalaran arsitektural, hasil evaluasi dapat
dibaca sebagai berikut. Sistem otomasi inspeksi kontainer memerlukan integrasi
antarkomponen, konsistensi identitas objek inspeksi, riwayat bukti, dan
pemisahan jalur komunikasi sesuai karakter data. Prinsip ini sejalan dengan
atribut kualitas ISO/IEC 25010, terutama \textit{functional suitability},
\textit{usability}, \textit{security}, dan \textit{maintainability}, serta
sejalan dengan prinsip pemisahan tanggung jawab dan \textit{loose coupling}
pada arsitektur sistem terdistribusi.

Berdasarkan premis tersebut, rancangan menetapkan empat keputusan utama:

\begin{enumerate}
  \item Komponen \textit{edge}, layanan API, aplikasi web, dan penyimpanan
  dipisahkan menurut tanggung jawabnya.
  \item Nomor kontainer dijadikan jangkar integrasi agar hasil dari beberapa
  proses inspeksi dapat dikaitkan pada entitas yang sama.
  \item Data terstruktur dan video langsung dipisahkan ke jalur yang berbeda
  agar masing-masing dapat dikelola sesuai fungsinya.
  \item Artefak visual dipisahkan dari data transaksional untuk menjaga ukuran
  data inti dan riwayat bukti.
\end{enumerate}

Hasil evaluasi pada bagian sebelumnya menunjukkan bahwa keempat keputusan itu
dapat direalisasikan pada sistem yang dibangun. Dengan demikian, validasi pada
tugas akhir ini mendukung proposisi bahwa rancangan arsitektur yang disusun
tidak berhenti pada tingkat konseptual, tetapi memiliki dasar logis dan bukti
realisasi sesuai kriteria pada Tabel~\ref{tbl:kriteria-keterpenuhan-evaluasi}
untuk memenuhi tujuan sistem pada ruang lingkup tugas akhir ini.

Bukti realisasi tersebut tidak hanya terlihat pada antarmuka, tetapi juga pada
struktur data inti dan penyimpanan artefak visual yang ditampilkan pada
Lampiran~\ref{sec:contoh-struktur-data}. Cuplikan entitas kontainer pada
Listing~\ref{lst:container-entity-sample} memperlihatkan bahwa hasil OCR,
\textit{defect scan}, \textit{illegal scan}, \textit{category scan},
\textit{manual review}, dan lampiran bukti visual memang terkonsolidasi pada
satu dokumen utama. Selain itu, Gambar~\ref{fig:evidence_r2_storage}
menunjukkan bahwa berkas visual disimpan pada penyimpanan objek Cloudflare R2,
sehingga pemisahan antara data transaksional dan artefak visual memiliki bukti
realisasi yang dapat diamati.

Pemeriksaan kode layanan API juga menunjukkan bahwa jalur utama integrasi
memang direalisasikan melalui keluarga \textit{endpoint} \nolinkurl{/api/edge}. Jalur
inti tersebut mencakup \textit{endpoint} identifikasi kontainer serta \textit{endpoint}
pemindaian kerusakan, barang berisiko, dan kategori muatan. Temuan ini
memperkuat kesimpulan bahwa arsitektur yang dibangun berpusat pada
entitas kontainer dan tidak bergantung pada jalur \textit{tracking} lama
sebagai pusat persistensi inspeksi.

Argumen validasi tersebut diringkas pada Tabel~\ref{tbl:argumen-validasi-arsitektur}.

\begin{table}[ht]
    \centering
    \footnotesize
    \caption[Argumen validasi arsitektur]{Argumen Validasi Arsitektur}
    \label{tbl:argumen-validasi-arsitektur}
    \setlength{\tabcolsep}{5pt}
    \renewcommand{\arraystretch}{1.12}
    \begin{tabularx}{\textwidth}{>{\raggedright\arraybackslash}X >{\raggedright\arraybackslash}X >{\raggedright\arraybackslash}X >{\raggedright\arraybackslash}X}
        \toprule
        \textbf{Premis Kebutuhan} & \textbf{Keputusan Arsitektur} & \textbf{Bukti Realisasi} & \textbf{Simpulan} \\
        \midrule
        Data inspeksi berasal dari beberapa proses dan harus tetap terhubung pada objek yang sama & Nomor kontainer dijadikan jangkar integrasi dan pusat agregasi data & Alur OCR membentuk atau memperbarui entitas kontainer, lalu hasil pemindaian lain dikaitkan ke entitas yang sama & Rancangan arsitektur data mendukung konsistensi konteks inspeksi pada skenario evaluasi \\
        Video langsung dan data transaksional memiliki karakteristik beban dan tujuan yang berbeda & Jalur video dipisahkan dari jalur data terstruktur & Video beranotasi disiarkan melalui \textit{WebSocket}, sedangkan hasil inspeksi dikirim ke layanan API & Pemisahan jalur komunikasi layak diterapkan pada ruang lingkup realisasi internal \\
        Hasil inspeksi harus dapat ditelusuri, diaudit, dan dibandingkan dengan konteks manifes & Data inti disimpan pada model kontainer dan manifes, sedangkan artefak visual dipisahkan ke penyimpanan objek & Koleksi \textit{containers} dan \textit{manifests} menyimpan data terstruktur, sedangkan bukti visual disimpan pada Cloudflare R2 & Rancangan penyimpanan mendukung riwayat kontainer dan pengelolaan bukti inspeksi pada artefak yang diuji \\
        Komponen sistem memiliki tanggung jawab berbeda dan harus tetap dapat berkembang tanpa saling membebani & Komponen \textit{edge}, API, web, dan penyimpanan dipisahkan menurut peran & Skenario integrasi menunjukkan komponen saling terhubung melalui kontrak layanan yang berbeda sesuai tanggung jawabnya & Pembagian lapisan arsitektur bersifat implementabel pada ruang lingkup tugas akhir ini \\
        \bottomrule
    \end{tabularx}
\end{table}

\FloatBarrier

\section{Keterbatasan Evaluasi}
\label{sec:keterbatasan-evaluasi}

Hasil evaluasi menunjukkan bahwa alur internal pada skenario yang diuji
konsisten dengan rancangan arsitektur. Namun, hasil tersebut belum dapat
digeneralisasi sebagai bukti operasi pelabuhan nyata.

Evaluasi pada tugas akhir ini berfokus pada integrasi sistem dan validasi alur
data. Estimasi kapasitas pada Subbab~\ref{subsec:estimasi-kapasitas-inspeksi}
memberikan gambaran kuantitatif awal, tetapi belum menggantikan uji beban dan
pengukuran performa operasional penuh seperti latensi menyeluruh, kapasitas
pemrosesan aktual, dan konsumsi sumber daya. Akurasi model deteksi dan OCR juga
belum diukur secara kuantitatif karena berada di luar fokus utama kontribusi
arsitektural. Karena itu, evaluasi lanjutan perlu diarahkan pada pengukuran
performa, uji beban, ketahanan operasional, dan validasi lapangan yang lebih
representatif.
